\section{模型检验与灵敏度分析}
\subsection{硬约束、复算与双运行}
solve v12 的约束审计显示：任务唯一性、到达、截止、实时开工、时延、GPU、IT 与设施功率均通过；Q3 和 Q4 的能源平衡、SOC、功率、互斥及终端状态最大绝对残差不超过 $2.274\times10^{-13}$。这些检验说明结果在给定容差与题面参数内自洽，但不提供最优性证书。

两次隔离运行使用相同输入副本、同一 code v15 与固定种子，自然耗时分别为 1687.58 s 和 1706.49 s。除记录自然耗时的 \texttt{method\_runtime.json} 外，48 个正式输出文件逐字节一致，证明给定环境下结果可复现；程序与结果的完整 SHA-256 保留在交付审计证据中。

\subsection{问题二参数灵敏度}
以平衡方案成本 $-421801727.45$ 元为基准，对购电价格整体缩放。系数为 0.8、0.9、1.1、1.2 时，成本分别为 $-421804612.25$、$-421803169.85$、$-421800285.05$、$-421798842.65$ 元，相对变化绝对值不超过 0.000684\%。图~\ref{fig:price} 表明当前净成本对购电价格整体缩放不敏感，原因是购电支出相对售电收入很小。
\begin{figure}[H]\centering
\includegraphics[width=.80\textwidth]{sensitivity_electricity_price.png}
\caption{问题二净成本的购电价格灵敏度}\label{fig:price}
\end{figure}

新能源水平系数为 0.8、0.9、1.1、1.2 时，成本分别为 $-3.039123$、$-3.811292$、$-4.464601$、$-4.576676$ 亿元，相对基准变化为 27.9490\%、9.6426\%、$-5.8460$\%、$-8.5030$\%。图~\ref{fig:renewable} 呈现明显非线性：新能源下降时购电缺口扩大，而上升时外送上限与弃电共同限制边际收益。
\begin{figure}[H]\centering
\includegraphics[width=.80\textwidth]{sensitivity_renewable_level.png}
\caption{问题二净成本的新能源水平灵敏度}\label{fig:renewable}
\end{figure}

\subsection{终端 SOC 与高循环策略解释}
终端 SOC 放宽后，Q3 成本最小候选仍为 $-458752156.98$ 元，说明其主要驱动是新能源与价格结构而非末端库存；但平衡与新能源策略的等效循环分别达到 1568.242 和 1568.882。题面未给循环寿命与退化成本，故高循环账面收益只能作为候选结果，不能直接外推为可部署收益。

Q4 三个新能源波动种子的峰值购电为 146.1294--298.6499 MW、碳排放为 464.9554--852.0600 tCO$_2$，显示单一随机轨迹不足以概括风险；当前三种子只提供方向性稳健性证据。
